\chapter*{Résumé}
\addcontentsline{toc}{chapter}{Résumé}
\markboth{RÉSUMÉ}{RÉSUMÉ}

Ce Projet de Fin d'Année (PFA) s'inscrit dans le cadre de la \textbf{conception d'une plateforme de gestion des bénéficiaires de l'Axe Social et Solidaire au sein de la Fondation OCP en intégrant des exigences de cybersécurité}. Face aux défis opérationnels liés au pilotage et au suivi d'un réseau national d'environ 1\,700 coopératives partenaires et de leurs bénéficiaires directs et indirects (femmes artisanes et agricultrices, jeunes porteurs de projets, personnes en situation de handicap, petits exploitants agricoles), la Fondation OCP a exprimé le besoin stratégique d'un système d'information centralisé, sécurisé et interopérable, capable de digitaliser les soumissions de rapports mensuels, de fiabiliser la collecte des données de terrain et de consolider les indicateurs d'impact ESG (Environnement, Social, Gouvernance) ainsi que les Objectifs de Développement Durable (ODD).

Pour répondre à la complexité de ce \textbf{système cible}, la phase de cadrage nécessitait un Cahier des Charges précis, exhaustif et rigoureusement structuré. Plutôt que de nous limiter à une démarche documentaire manuelle sujette aux omissions, nous avons conçu et développé \textbf{Secure Requirements Studio (SRS)}, une solution logicielle innovante dédiée à l'ingénierie sécurisée des exigences et à l'automatisation des spécifications. SRS s'appuie sur une ontologie métier de 39 concepts clés interconnectés, propose un moteur d'interview à double mode (questionnaire déterministe guidé et entretien intelligent assisté par IA), intègre nativement la modélisation des menaces (STRIDE/DREAD) et la matrice des privilèges (RBAC), et audite la maturité du cadrage via plus de 100 règles de validation formelles.

Cette démarche a abouti à la génération du document officiel de référence {\ttfamily Cahier\allowbreak\_\allowbreak des\allowbreak\_\allowbreak Charges\allowbreak\_\allowbreak FOCP\allowbreak\_\allowbreak V3.pdf}. Enfin, afin de valider concrètement les spécifications produites, nous avons réalisé le \textbf{MVP (Minimum Viable Product)} de la plateforme de gestion des bénéficiaires, concrétisant ainsi la transition fluide entre l'ingénierie des exigences et le prototypage applicatif fonctionnel (tableaux de bord, annuaire des coopératives, cartographie interactive, reporting mensuel et gestion documentaire).

\vspace{0.8cm}
\noindent\textbf{Mots-clés :} Gestion des bénéficiaires, Fondation OCP, Axe Social et Solidaire, Secure Requirements Studio (SRS), Ingénierie des exigences, Cybersécurité, STRIDE, RBAC, Cahier des Charges, MVP, Next.js, FastAPI.

\vfill
\newpage

\chapter*{Abstract}
\addcontentsline{toc}{chapter}{Abstract}
\markboth{ABSTRACT}{ABSTRACT}

This Final Year Project (PFA) focuses on the \textbf{design of a beneficiary management platform for the Social and Solidarity Axis within Fondation OCP while integrating cybersecurity requirements}. Facing the operational challenges of monitoring and steering a nationwide network of approximately 1,700 partner cooperatives and their direct and indirect beneficiaries (women artisans and farmers, young entrepreneurs, people with disabilities, smallholder farmers), Fondation OCP identified the strategic need for a centralized, secure, and data-driven information system capable of streamlining monthly reporting workflows, securing field data collection, and consolidating ESG (Environmental, Social, Governance) impact indicators as well as UN Sustainable Development Goals (SDGs).

To tackle the operational complexity of this \textbf{target system}, the requirements engineering phase demanded a comprehensive, precise, and highly qualified specification document. Rather than relying on manual drafting methods prone to critical omissions, we designed and developed \textbf{Secure Requirements Studio (SRS)}, a specialized software platform dedicated to automated and security-driven requirements engineering. SRS relies on a 39-concept ontological knowledge graph, features a dual-mode elicitation engine (deterministic guided questionnaire and AI-assisted interview), embeds automated threat modeling (STRIDE/DREAD) and role-based access control (RBAC) matrices, and audits specification maturity against over 100 formal validation rules.

This structured workflow led to the generation of the official reference specification deliverable {\ttfamily Cahier\allowbreak\_\allowbreak des\allowbreak\_\allowbreak Charges\allowbreak\_\allowbreak FOCP\allowbreak\_\allowbreak V3.pdf}. Furthermore, to provide a tangible proof-of-concept and validate the generated specifications, we developed the \textbf{MVP (Minimum Viable Product)} of the beneficiary management platform, demonstrating a seamless transition from formal requirements to a concrete functional prototype (dashboards, cooperative registry, interactive nationwide map, monthly reporting workflows, and document management).

\vspace{0.8cm}
\noindent\textbf{Keywords:} Beneficiary Management, Fondation OCP, Social and Solidarity Axis, Secure Requirements Studio (SRS), Requirements Engineering, Cybersecurity, STRIDE, RBAC, Specification Deliverable, MVP Prototype, Next.js, FastAPI.

\vfill
\newpage
